\documentclass[11pt,a4paper]{article}
\usepackage[utf8]{inputenc}
\usepackage[T1]{fontenc}
\usepackage[italian]{babel}
\usepackage{geometry}
\usepackage{booktabs}
\usepackage{hyperref}
\usepackage{enumitem}
\geometry{margin=2.4cm}

\title{Confronto tra PQ-BLE v0.7 e v1.0\\
\large Architettura, riconnessione, overhead e risultati sperimentali}
\author{Alessio Bonora}
\date{Settembre 2026}

\begin{document}
\maketitle
\tableofcontents
\newpage

\section{Obiettivo}
Questo documento descrive le due implementazioni principali del progetto
PQ-BLE-HANDSHAKE e ne confronta architettura, gestione della riconnessione,
overhead, prestazioni, risorse e proprietà di sicurezza validate.

Le due versioni rappresentano scelte architetturali differenti:
\begin{itemize}[nosep]
\item \textbf{v0.7}: protocollo ibrido interamente applicativo, basato su ML-KEM-768 e P-256 ECDH;
\item \textbf{v1.0}: architettura a livelli, con BLE Security Mode 1 Level 4 per autenticazione/protezione del link e ML-KEM-768 per il key establishment applicativo post-quantum.
\end{itemize}

La v1.0 termina con CP5. La successiva Post-v1.0 Comparative Evaluation
(EVAL-A--EVAL-E) confronta le due baseline congelate.

\section{v0.7: handshake ibrido applicativo}
Nel profilo v0.7 BLE SMP è disabilitato. La sicurezza è costruita sopra GATT a livello applicativo.

\begin{verbatim}
SS_MLKEM + SS_ECDH
        |
u16be(32) || SS_MLKEM || u16be(32) || SS_ECDH
        |
HKDF-SHA256
        |
SAS + FINISHED + AES-256-GCM
\end{verbatim}

\subsection{Overhead del data-plane v0.7}
\begin{verbatim}
sequence number   8 B
message type      1 B
IV                12 B
GCM tag           16 B
----------------------
overhead          37 B
\end{verbatim}

Nel workload hardware congelato: 6 B di plaintext, 43 B di frame.

\section{Session resumption storica e v0.7 finale}
Il PoC Python iniziale disponeva di un \texttt{SessionStore} persistente.
La stessa session key poteva essere riutilizzata fino al primo limite tra
\textbf{24 ore} e \textbf{100 resume riusciti}. Timeout e resume falliti non
consumavano il contatore.

Il resume richiedeva circa:
\begin{verbatim}
21 B request + 9 B response = 30 B applicativi
\end{verbatim}

Questa ottimizzazione evitava il full handshake, ma riutilizzava la stessa session key.

\subsection{Distinzione fondamentale}
La session resumption del PoC storico \textbf{non è una funzionalità della v0.7
hardware finale}. Nel profilo v0.7 usato nella comparazione BLE SMP è
disabilitato, non esiste un bond BLE e non è stata completata una
session-resumption persistente sul DK.

\section{v1.0: BLE L4 + ML-KEM}
\begin{verbatim}
BLE SMP Security Mode 1 Level 4
        |
ML-KEM-768
        |
SS_MLKEM
        |
transcript + HKDF-SHA256
        |
FINISHED_C / FINISHED_P
        |
K_APP_C2P / K_APP_P2C
        |
AES-256-GCM applicativo
\end{verbatim}

Il materiale SMP non viene esportato né combinato con \texttt{SS\_MLKEM}.

\section{Riconnessione nella v1.0}
La proprietà centrale è:
\begin{verbatim}
bond BLE persistente != sessione applicativa persistente
\end{verbatim}

Una riconnessione bonded ripristina L4 senza ripetere la Numeric Comparison, ma crea una nuova sessione:
\begin{verbatim}
bonded L4
 -> nuovo ML-KEM
 -> nuovo CP3
 -> nuove chiavi applicative
 -> counter CP4 da zero
\end{verbatim}

Il bonding riduce quindi il costo di associazione/autenticazione BLE nelle
connessioni successive, ma non elimina il nuovo key establishment ML-KEM applicativo.

\section{Overhead del data-plane v1.0}
\begin{verbatim}
PQV1 header       8 B
seq               8 B
msg_type          1 B
plaintext_len     2 B
ciphertext        N B
tag              16 B
\end{verbatim}

L'overhead applicativo è \textbf{35 B}. L'IV non viene trasmesso.

\section{Confronto dell'overhead per messaggio}
A parità di plaintext $N$:
\begin{verbatim}
v0.7 = N + 37 B
v1.0 = N + 35 B
\end{verbatim}

\begin{center}
\begin{tabular}{rrrrr}
\toprule
Plaintext & v0.7 & v1.0 & Ov. v0.7 & Ov. v1.0 \\
\midrule
16 B   & 53 B   & 51 B   & 231.3\% & 218.8\% \\
64 B   & 101 B  & 99 B   & 57.8\%  & 54.7\%  \\
256 B  & 293 B  & 291 B  & 14.5\%  & 13.7\%  \\
512 B  & 549 B  & 547 B  & 7.2\%   & 6.8\%   \\
1024 B & 1061 B & 1059 B & 3.6\%   & 3.4\%   \\
\bottomrule
\end{tabular}
\end{center}

Il fatto che v1.0 risparmi 2 B nel frame applicativo non implica automaticamente
meno overhead radio totale, perché v1.0 usa anche cifratura/autenticazione BLE Link Layer.

\section{Overhead dell'handshake e della riconnessione}
La v0.7 gestisce ML-KEM, P-256 applicativo, transcript ibrido, SAS e FINISHED.

La v1.0 elimina P-256 ECDH e SAS custom dal protocollo applicativo, ma in cold mode
esegue SMP L4 con Numeric Comparison. In bonded mode la Numeric Comparison non viene
ripetuta, ma ML-KEM, CP3 e FINISHED vengono eseguiti di nuovo.

\section{Risultati sperimentali}
\begin{verbatim}
v0.7 hybrid: n = 30
v1.0 bonded: n = 30
v1.0 cold:   n = 5
\end{verbatim}

Tempo medio \texttt{secure\_machine\_ms}:
\begin{verbatim}
v0.7        3584.773 ms
v1.0 bonded 5446.180 ms
\end{verbatim}

ML-KEM encapsulation lato Central:
\begin{verbatim}
v0.7        0.653 ms
v1.0 bonded 0.638 ms
\end{verbatim}

La differenza end-to-end non deve quindi essere attribuita a ML-KEM.

\section{Risorse}
\begin{center}
\begin{tabular}{lrrr}
\toprule
Risorsa & v0.7 & v1.0 & Delta \\
\midrule
FLASH & 225652 B & 279428 B & +53776 B (+23.8\%) \\
RAM & 106048 B & 109464 B & +3416 B (+3.2\%) \\
Crypto worker stack & 28672 B & 28672 B & 0 \\
Peak osservato & 24264 B & 24264 B & 0 \\
\bottomrule
\end{tabular}
\end{center}

\section{Visibilità del traffico}
La cattura v0.7 ha osservato 716 packet in 8.145325 s, con 95 ATT e 0 SMP.
La v1.0 cold ha osservato 512 packet in 7.605352 s, con 50 ATT e 9 SMP.
La v1.0 bonded ha osservato 37 packet in 0.945367 s, con 0 ATT decodificato e 1 SMP.

Dopo l'attivazione della cifratura BLE, senza LTK il traffico ATT v1.0 non è più decodificabile.

\section{Validazione di sicurezza}
\begin{verbatim}
v0.7: 7/7 PASS
v1.0: 2/2 PASS
totale: 9/9 PASS
\end{verbatim}

\section{Conclusione}
v0.7 è più leggera e più veloce nell'implementazione misurata, ma concentra
l'autenticazione e il key agreement classico/PQ nel protocollo applicativo.

v1.0 usa invece una separazione più modulare: BLE L4 fornisce autenticazione e
protezione del link, mentre ML-KEM fornisce il key establishment applicativo
post-quantum.

Il resume 24 ore / 100 utilizzi apparteneva al PoC Python storico, non alla
v0.7 hardware finale. I dati attuali non dimostrano che v1.0 abbia meno byte
radio totali rispetto a v0.7.

\end{document}
